\begin{table}[htbp]
\centering
\caption{Cultural correlates of mean fertility, cross-section of 14 countries}
\label{tab:crosssection}
\small
\begin{tabular}{lcccc}
\toprule
 & (1) & (2) & (3) & (4) \\
\midrule
Muslim share (\%)      & $+0.0130$ &          & $+0.0120$ & $+0.0048$ \\
                       & (0.0032)  &          & (0.0039)  & (0.0209) \\
Female SMAM (years)    &           & $-0.0897$ & $-0.0188$ &          \\
                       &           & (0.0390)  & (0.0238)  &          \\
Central Asia           &           &           &           & $+0.941$ \\
                       &           &           &           & (1.868) \\
\midrule
$R^2$                  & 0.774 & 0.323 & 0.783 & 0.911 \\
Observations           & 14 & 14 & 14 & 14 \\
\bottomrule
\end{tabular}

\begin{minipage}{\textwidth}\footnotesize
\vspace{4pt}
\emph{Note:} Dependent variable is mean total fertility rate, 2000--2023. HC3 standard
errors in parentheses, with small-sample $t$ inference. Column (4) is the inseparability
test: $\mathrm{corr}(\text{Central Asia}, \text{Muslim share}) = +0.831$, and neither
variable is distinguishable from zero when both are included
(Central Asia: $p = 0.624$; Muslim share: $p = 0.822$) even though the specification fits
better than any other column.
\end{minipage}
\end{table}
